\section{Temperature, sampling y oracle approximation}

La sampling strategy determina de qué región de la model distribution se toman los candidatos. La \term{temperature} (temperatura) se aplica a los logits $z_i$ así:
\[
p_T(i)=\frac{\exp(z_i/T)}{\sum_j\exp(z_j/T)}.
\]
Con $T<1$ la distribución se vuelve más aguda y favorece los token de alta probabilidad; con $T>1$ aumenta la diversidad, pero también los errores sintácticos o semánticos. La temperatura óptima depende de la metric y de $k$.

\subsection{Calibrar la temperature para pass@k}

Esta sección vincula primero los parámetros de decoding con el presupuesto de evaluación. La temperatura no debe fijarse por costumbre, sino calibrarse según $k$, el verifier y el presupuesto de latencia. Una temperatura baja concentra los candidatos en la mode, lo que conviene cuando se muestra una sola salida; una temperatura alta amplía la search coverage, pero solo unas tests confiables pueden convertir la diversity en corrección final.

\lecturefigure{slide-39-temperature-calibration.jpg}{Una temperatura baja favorece pass@1; una temperatura más alta puede aportar candidatos diversos cuando $k$ es grande (large-k).}{Video oficial 41:26.}

\begin{knowledgebox}{Lectura de la figura: cada curva corresponde a un presupuesto de búsqueda}
Cuando $k=1$, una temperatura demasiado alta dispersa la masa de probabilidad hacia errores de baja probabilidad; cuando $k$ es grande, una diversity moderada evita repetir el mismo patrón de error. La temperatura óptima no es una constante del modelo: debe calibrarse en validation tasks según el generation budget real, y hay que evitar el benchmark overfitting.
\end{knowledgebox}

\subsection{La «efectividad irrazonable» del sampling}

A continuación se observa la curva de beneficio al aumentar el candidate budget, y la benchmark gain se pone en una misma decisión junto con los costos de inferencia, ejecución y selección. El desplazamiento horizontal de la curva significa más samples; la subida vertical significa que en más problemas aparece al menos un candidato correcto. Entre ambos no hay un aumento gratuito de capacidad.

\lecturefigure{slide-40-sampling-effectiveness.jpg}{Al aumentar los samples, pass@k puede superar con mucho el desempeño de una sola salida.}{Video oficial 42:19.}

\begin{knowledgebox}{Lectura de la figura: la mejora proviene de la search, no necesariamente de una mejor muestra individual}
La curva muestra que el mismo modelo, con un candidate budget mayor, resuelve más problemas. Un sistema real además tiene que pagar los costos de inference, sandbox, test y ranking. Si las tests son incompletas, el selector puede favorecer el exploit; si varios candidatos están muy correlacionados, el número efectivo de muestras es menor que el $k$ nominal.
\end{knowledgebox}

\teachervoice{El ponente llama al sampling «unreasonably effective». La nota de clase no invita a muestrear más a ciegas, sino que recuerda que la calidad del modelo es toda una conditional distribution; cuando un producto muestra una sola completion, el decoding y el verifier determinan juntos la capacidad visible.}

\subsection{Aproximar el oracle}

Después se introduce el límite superior de la selección. Aquí el \term{oracle} es un selector hipotético: siempre que entre los $k$ candidatos exista un programa correcto, puede identificarlo y devolverlo. Las unit tests son una oracle approximation finita: cuanto más completa su cobertura, más confiable es la selección. Esta abstracción separa «el modelo no generó la respuesta correcta» de «el sistema no reconoció la respuesta correcta» como dos cuellos de botella medibles.

\lecturefigure{slide-41-oracle-sampling.jpg}{El sampling más un verifier aproxima el oracle que «elige el programa correcto del conjunto de candidatos».}{Video oficial 43:17.}

\begin{knowledgebox}{Lectura de la figura: comparar la distancia entre la model curve y la oracle curve}
Si el pass@k de tipo oracle es alto pero el del selector real es bajo, el cuello de botella está en la verification/ranking; si ambos son bajos, el cuello de botella está en la model distribution. Los equipos de ingeniería deben invertir por separado en generation, tests, static analysis y reranking, y no atribuir toda la diferencia a los parámetros del modelo.
\end{knowledgebox}

\begin{warningbox}{Las pruebas pueden satisfacerse, pero también aprovecharse con trampas}
Unas tests finitas permiten overfit, hard-code o efectos secundarios no cubiertos. El código generado debe ejecutarse en un sandbox, revisando timeout, memory, network, filesystem y dependency. El oracle es un límite superior abstracto, no una garantía de seguridad en la práctica.
\end{warningbox}

\subsection{Codex-S: la señal de supervisión sigue dando forma}

Codex-S aplica supervised fine-tuning sobre un code-pretrained model, con stand-alone Python functions y sus solutions correctas. Esto modifica la prompt-to-solution distribution y la ajusta mejor a las tareas de tipo HumanEval (HumanEval-style tasks).

\lecturefigure{slide-42-codex-s-training.jpg}{Codex-S usa datos de funciones supervisados para alinear aún más la completion de tipo HumanEval (HumanEval-style).}{Video oficial 44:10; Chen et al. 2021.}

\begin{knowledgebox}{Lectura de la figura: el fine-tuning mejora el ajuste a la tarea (task fit), pero también puede estrechar la distribución}
Los supervised examples enseñan explícitamente al modelo a pasar del docstring al function body, lo que eleva la benchmark performance. Hay que revisar si se sacrifican otros lenguajes, el repository context o la diversity, y evitar la train-test similarity. La fine-tuning gain no significa que el base model no importe, sino que el objective se acerca más a la evaluation.
\end{knowledgebox}

\subsection{Main figure: escala del modelo, Codex y sampling budget}

Esta sección integra model family, scale y sampling budget, y distingue la pretraining gain, la code fine-tuning gain y la search gain, para no mezclar los tres efectos en una sola cifra.

\lecturefigure{slide-43-main-figure.jpg}{La main figure de Codex compara en conjunto model family, escala y pass@k.}{Video oficial 45:52; Chen et al. 2021.}

\begin{knowledgebox}{Lectura de la figura: primero la leyenda, después las pendientes y las distancias}
El eje horizontal o las curvas suelen codificar el model size o el sampling budget; los colores distinguen GPT, Codex y Codex-S. Compare primero la fine-tuning gain con el mismo $k$, y luego la search gain de un mismo modelo de pass@1 a pass@100. No tome pass@100 como la experiencia interactiva de una sola completion, ni ignore el costo de generar 100 candidatos.
\end{knowledgebox}

\subsection{Resumen de la sección}

La temperature controla el equilibrio calidad–diversidad, pass@k incorpora el presupuesto de candidatos a la evaluación, las tests aproximan el oracle y Codex-S modifica la distribución mediante supervised data. La capacidad de un sistema de código proviene de tres elementos: model, decoding y verifier.
